\documentclass[12pt, oneside]{book}
\usepackage{amsmath, amssymb, amssymb, amsthm, mathtools, graphicx, hyperref, titling}
\usepackage{tikz, caption, subcaption, enumitem, polyglossia}
\usepackage{tikz-3dplot, float, etoolbox, enumitem, tcolorbox}
\usepackage[top=25mm, bottom=25mm, left=20mm, right=20mm]{geometry}
%\usepackage{showframe}

\makeatletter
\patchcmd{\@makechapterhead}{\vspace*{50\p@}}{}{}{}% Removes space above \chapter head
\patchcmd{\@makeschapterhead}{\vspace*{50\p@}}{}{}{}% Removes space above \chapter* head
\makeatother

\newtheorem{theorem}{Theorem}[chapter]
\newtheorem{lemma}[theorem]{Lemma}
\newtheorem{corollary}[theorem]{Corollary}
\newtheorem{proposition}[theorem]{Proposition}

\predate{}
\postdate{}
\date{}
\title{MT3144 - Ordinary Differential Equations}
\author{Nachiketa Kulkarni}
\pagenumbering{gobble}
\setmainfont{Comic Sans MS}

\begin{document}
\maketitle
\tableofcontents

\mainmatter
\pagenumbering{arabic}
\chapter{Introduction}
Ordinary differential equations (ODEs) are a mathematical model for Dynamics.
Examples:
\begin{itemize}
	\item Positive Feedback: \(\frac{dx}{dt} = kx, k > 0\)
	\item Negative Feedback: \(\frac{dx}{dt} = -kx, k > 0\)
	\item Hooke's Law: \(\frac{d^2x}{dt^2} = -kx, k > 0\)
	\item Law of Gravity: \(\frac{d^2x}{dt^2} = -\frac{kx}{|x|^3}\)
	\item Lotka-Volterra Equations: \(\frac{dx}{dt} = kx - xy, \frac{dy}{dt} = -ky + xy\)
\end{itemize}
\section{Most Important ODE - \(x'(t) = f(t)\)}
Domain of \(f, u\) is \(\mathbb{R}\) (or an interval of \(\mathbb{R}\)).
We will consider the initial value problem (IVP) for this ODE:
\begin{align*}
	x'(t) &= f(t) \\
	x(t_0) &= x_0
\end{align*}
\begin{theorem}
	\(u: I \rightarrow \mathbb{R}\), \(u\) is differentiable on \(I\) and \(u'(t) = f(t)\) for all \(t \in I\).
	Then there exists a unique solution \(u\) to the IVP on \(I\):
	\begin{equation}
		u(t) = x_0 + \int_{t_0}^{t} f(s) ds
	\end{equation}
\end{theorem}
\begin{proof}
	\textbf{Existence:} The existance of the solution is guaranteed by the Fundamental Theorem of Calculus.\\
	\textbf{Uniqueness:} Suppose \(v: I \rightarrow \mathbb{R}\) is another solution to the IVP.
	Let \(w(t) = u(t) - v(t)\).
	Then \(w'(t) = u'(t) - v'(t) = f(t) - f(t) = 0\).
	Therefore, \(w(t) = C\) for some constant \(C\).
	Since \(w(t_0) = u(t_0) - v(t_0) = x_0 - x_0 = 0\), we have \(C = 0\).
	Thus, \(w(t) = 0\) for all \(t \in I\), which implies \(u(t) = v(t)\) for all \(t \in I\).
	Hence, the solution is unique.
\end{proof}
\section{Quadratic First Order ODEs}
Consider the IVP:
\begin{align*}
	(x'(t))^2 &= t \\
	x(0) &= 1
\end{align*}

The solution to this IVP is done by taking multiple cases:
\begin{itemize}
	\item Case 1: \(x'(t) = \sqrt{t}\)
	\begin{align*}
		x(t) &= \int \sqrt{t} dt \\
		&= \frac{2}{3} t^{3/2} + C
	\end{align*}
	Using the initial condition \(x(0) = 1\), we find \(C = 1\). Therefore, the solution for Case 1 is:
	\begin{equation}
		x(t) = \frac{2}{3} t^{3/2} + 1
	\end{equation}
	\item Case 2: \(x'(t) = -\sqrt{t}\)
	\begin{align*}
		x(t) &= \int -\sqrt{t} dt \\
		&= -\frac{2}{3} t^{3/2} + C
	\end{align*}
	Using the initial condition \(x(0) = 1\), we find \(C = 1\). Therefore, the solution for Case 2 is:
	\begin{equation}
		x(t) = -\frac{2}{3} t^{3/2} + 1
	\end{equation}
\end{itemize}
\end{document}